% ============================================
% LaTeX 表格模板 (解决浮动体问题)
% ============================================

% ---------- 方法1: 非浮动表格（不会抢编号）----------
% 适用于变量定义表、简单统计表
\begin{center}
\captionof{table}{表格标题}
\small
\begin{tabular}{lccccc}
\toprule
变量 & 均值 & 标准差 & 最小值 & 最大值 & 观测值 \\
\midrule
变量1 & 0.123 & 0.456 & -2.34 & 3.45 & 1000 \\
变量2 & 1.234 & 2.345 & 0.00 & 10.0 & 1000 \\
\bottomrule
\end{tabular}
\label{tab:descriptive}
\end{center}

% ---------- 方法2: 浮动体表格（esttab输出用）----------
% 适用于回归结果表
\begin{table}[htbp]
\centering
\caption{回归结果表}
\footnotesize
\input{../analysis/output/regression_table.tex}
\label{tab:baseline}
\end{table}

% ---------- 方法3: 宽表优化（6列以上）----------
\begin{table}[htbp]
\centering
\caption{宽表示例（自动调整字号）}
\scriptsize  % 对于7列以上使用 \scriptsize
\setlength{\tabcolsep}{4pt}  % 减少列间距
\begin{tabular}{lccccccc}
\toprule
 & (1) & (2) & (3) & (4) & (5) & (6) & (7) \\
变量 & 系数 & 系数 & 系数 & 系数 & 系数 & 系数 & 系数 \\
\midrule
X & 0.123*** & 0.123** & 0.123* & 0.123 & 0.123*** & 0.123** & 0.123* \\
  & (0.012) & (0.023) & (0.067) & (0.123) & (0.012) & (0.023) & (0.067) \\
\bottomrule
\multicolumn{8}{l}{\footnotesize 注：括号内为聚类稳健标准误；* p<0.1, ** p<0.05, *** p<0.01}
\end{tabular}
\label{tab:wide_table}
\end{table}

% ---------- 重要：浮动体位置控制 ----------
% 在章节末尾使用 \FloatBarrier 强制输出所有待排表格
% 需要在导言区添加: \usepackage{placeins}
%
% 使用示例：
% \section{实证结果}
% 正文内容...
% \FloatBarrier  % 强制输出本节所有浮动体
% \section{下一章节}
